% Part: incompleteness
% Chapter: arithmetization-syntax
% Section: coding-symbols

\documentclass[../../../include/open-logic-section]{subfiles}

\begin{document}

\olfileid{inc}{art}{cod}
\olsection{चिन्हांचे सांकेतीकरण}

प्रथम-क्रम तर्कशास्त्राच्या मूलभूत भाषा~$\Lang L$ मध्ये ही चिन्हे वापरली जातात:
\[
\lfalse \quad \lnot \quad \lor \quad \land \quad \lif \quad \lforall
\quad \lexists \quad \eq \quad ( \quad ) \quad ,
\]
यांबरोबर चरे व !!{constant}s यांचे !!{enumerable} संच, तसेच कोणत्याही
स्थानसंख्येची !!{function}s व !!{predicate}s यांचे !!{enumerable} संचही
वापरले जातात. या प्रत्येक चिन्हाला \emph{संकेतांक} म्हणून एक संख्या
देता येते: प्रत्येक चिन्हाचा संकेतांक वेगळा असेल आणि दोन भिन्न
चिन्हांना एकच संख्या दिली जाणार नाही. सर्व चिन्हांचा संच !!{enumerable}
असल्यामुळे त्याची नैसर्गिक संख्यांच्या संचाशी !!a{bijection} आहे;
म्हणून हे शक्य आहे. परंतु संकेतांकावरून चिन्ह आणि त्याविषयीची काही
माहिती (उदा. !!a{function} ची स्थानसंख्या) संगणनक्षम रीतीने परत
मिळाली पाहिजे. असे करण्याचे अनेक मार्ग आहेत. पुढील मार्गात आदिम
पुनरावर्ती फलने वापरली आहेत. ($\tuple{n_0, \dots, n_k}$ ही
$n_0$, \dots, $n_k$ या संख्यांच्या अनुक्रमाला सांकेतिक करणारी संख्या
आहे, हे आठवा.)

\begin{defn}
$s$ हे~$\Lang L$ मधील चिन्ह असेल, तर त्याचा \emph{चिन्हसंकेतांक}~$\scode s$
पुढीलप्रमाणे परिभाषित करू:
\begin{enumerate}
\item $s$ हे तार्किक चिन्हांपैकी असेल, तर पुढील सारणीनुसार $\scode s$ ठरतो:
\[
\begin{array}{cccccccccc}
  \lfalse & \lnot & \lor & \land & \lif & \lforall \\
\tuple{0, 0} & \tuple{0, 1} & \tuple{0, 2} & \tuple{0, 3} &
\tuple{0, 4} & \tuple{0, 5} \\
\lexists & \eq & ( & ) & ,\\
\tuple{0, 6} & \tuple{0, 7} &
\tuple{0, 8} & \tuple{0, 9} & \tuple{0, 10}
\end{array}
\]
\item $s$ हा $i$-वा चर $\Obj v_i$ असेल, तर $\scode s = \tuple{1, i}$.
\item $s$ हे $i$-वे !!{constant}~$\Obj c_i$ असेल, तर
  $\scode s = \tuple{2, i}$.
\item $s$ हे $i$-वे $n$-स्थानी !!{function}~$\Obj f_i^n$ असेल, तर
  $\scode s = \tuple{3, n, i}$.
\item $s$ हे $i$-वे $n$-स्थानी !!{predicate}~$\Obj P_i^n$ असेल, तर
  $\scode s = \tuple{4, n, i}$.
\end{enumerate}
\end{defn}

\begin{prop}
पुढील संबंध आदिम पुनरावर्ती आहेत:
\begin{enumerate}
\item $\fn{Fn}(x, n)$ तेव्हा आणि केवळ तेव्हाच जेव्हा $x$ हा एखाद्या~$i$ साठी
  $\Obj f^n_i$ चा संकेतांक असतो; म्हणजेच $x$ हा $n$-स्थानी
  !!{function} चा संकेतांक असतो.
\item $\fn{Pred}(x, n)$ तेव्हा आणि केवळ तेव्हाच जेव्हा $x$ हा एखाद्या~$i$ साठी
  $\Obj P^n_i$ चा संकेतांक असतो, किंवा $x$ हा $\eq$ चा संकेतांक
  असतो आणि $n = 2$ असते; म्हणजेच $x$ हा $n$-स्थानी !!{predicate}
  चा संकेतांक असतो.
\end{enumerate}
\end{prop}

\begin{defn}
$s_0, \dots, s_{n-1}$ हा चिन्हांचा अनुक्रम असेल, तर त्याची
\emph{गोडेल संख्या} $\tuple{\scode{s_0}, \dots, \scode{s_{n-1}}}$ आहे.
\end{defn}

\begin{explain}
\emph{संकेतांक} आणि \emph{गोडेल संख्या} वेगवेगळ्या गोष्टी आहेत,
हे लक्षात घ्या. उदाहरणार्थ, $\Obj v_5$ या चराचा संकेतांक
$\scode{\Obj v_5} = \tuple{1, 5} = 2^2\cdot 3^6$ आहे.
पण पद म्हणून पाहिले असता~$\Obj v_5$ हाही लांबी~$1$ असलेला
चिन्हांचा अनुक्रम आहे. त्यामुळे $\Obj v_5$ या \emph{पदाची}
\emph{गोडेल संख्या}~$\Gn{\Obj v_5}$ ही
$\tuple{\scode{\Obj v_5}} = 2^{\scode{\Obj v_5} + 1} =
2^{2^2\cdot 3^6 + 1}$ आहे.
\end{explain}

\begin{ex}
$k_0$, \dots, $k_{n-1}$ हा संख्यांचा अनुक्रम असेल, तर
अविभाज्य संख्यांच्या घातांनी केलेल्या सांकेतीकरणात
$\tuple{k_0, \dots, k_{n-1}}$ हा त्याचा संकेतांक
\[
2^{k_0+1}\cdot3^{k_1+1}\cdot \dots \cdot p_{n-1}^{k_{n-1}+1},
\]
आहे, हे आठवा. येथे $p_i$ ही $i$-वी अविभाज्य संख्या आहे
($p_0 = 2$ पासून सुरुवात होते). म्हणून, उदाहरणार्थ,
$\eq[\Obj v_0][\Obj 0]$ हे सूत्र, किंवा अधिक स्पष्टपणे
${\eq}(\Obj v_0,\Obj c_0)$, याची गोडेल संख्या
\[
\tuple{\scode{\eq},\scode{(},\scode{\Obj v_0},\scode{,}, \scode{\Obj
    c_0},\scode{)}}.
\]
आहे. येथे $\scode{\eq}$ हा $\tuple{0,7} = 2^{0+1}\cdot
3^{7+1}$ आहे; $\scode{\Obj v_0}$ हा $\tuple{1,0} = 2^{1+1}\cdot3^{0+1}$
आहे; इत्यादी. म्हणून $\Gn{=(\Obj v_0,\Obj c_0)}$ ही संख्या
\begin{multline*}
2^{\scode{=} + 1}\cdot 3^{\scode{(}+1}\cdot 5^{\scode{\Obj v_0}+1}
\cdot 7^{\scode{,} + 1} \cdot 11^{\scode{\Obj c_0}+1} \cdot
13^{\scode{)}+1} = \\
2^{2^1\cdot 3^8 + 1}\cdot 3^{2^1\cdot 3^9+1}\cdot 5^{2^2\cdot 3^1+1}
\cdot 7^{2^1\cdot 3^{11} + 1} \cdot 11^{2^3\cdot3^1+1} \cdot
13^{2^1\cdot3^{10}+1} = \\
2^{13\,123}\cdot 3^{39\,367}\cdot 5^{13}\cdot 7^{354\,295}\cdot11^{25}\cdot13^{118\,099}.
\end{multline*}
\end{ex}

\end{document}
